% Unit 055: English translation of chapter4.tex, lines 987--1313.
% Source: Methods of Algebra, Volume 2, commit
% 9a5803ff77dd3257484cb177f851a73770a59dd3 (CC BY 4.0).
% stable-unit-id: o014.aljabr2.chapter4.morphisms-extensions
% Coverage: complete Section 4.5 on morphisms in the derived category,
% amplitude, the Way-out Lemma, Ext functors, and Yoneda's description by
% n-extensions; ends immediately before Section 4.6 at line 1314.

% segment-id: o014.li2.chapter4.morphisms-extensions.h001
\section{Morphisms and Extensions}\label{sec:Hom-Ext}
\hypertarget{o014.aljabr2.chapter4.morphisms-extensions}{ }

% segment-id: o014.li2.chapter4.morphisms-extensions.p001
Throughout this section, $\mathcal{A}$ is an Abelian category. Without
further comment, we identify
$\Obj(\cate{C}(\mathcal{A}))$, $\Obj(\cate{K}(\mathcal{A}))$, and
$\Obj(\cate{D}(\mathcal{A}))$; when necessary, we distinguish the morphism
sets in these categories by writing
$\Hom_{\cate{D}(\mathcal{A})}$, $\Hom_{\cate{K}(\mathcal{A})}$,
$\Hom_{\cate{C}(\mathcal{A})}$, and $\Hom_{\mathcal{A}}$. Readers not yet
familiar with K-injective and K-projective complexes may skip the statements
concerning them.

% segment-id: o014.li2.chapter4.morphisms-extensions.pr001
\begin{proposition}\label{prop:Hom-D-injective}
	Suppose $X \in \Obj(\cate{C}^+(\mathcal{A}))$ consists of injective
	objects, or, more generally, suppose
	$X \in \Obj(\cate{C}(\mathcal{A}))$ is a K-injective complex as in
	Definition~\ref{def:K-resolutions}. Then
	$\Hom_{\cate{K}(\mathcal{A})}(\cdot, X) \simeq
	\Hom_{\cate{D}(\mathcal{A})}(\cdot, X)$.

	Dually, suppose $X \in \Obj(\cate{C}^-(\mathcal{A}))$ consists of
	projective objects, or suppose
	$X \in \Obj(\cate{C}(\mathcal{A}))$ is a K-projective complex. Then
	$\Hom_{\cate{K}(\mathcal{A})}(X, \cdot) \simeq
	\Hom_{\cate{D}(\mathcal{A})}(X, \cdot)$.
\end{proposition}

% segment-id: o014.li2.chapter4.morphisms-extensions.pf001
\begin{proof}
	By duality, it suffices to treat
	$\Hom_{\cate{D}(\mathcal{A})}(\cdot, X)$. Lemma~\ref{prop:MXY-equiv}
	gives
% segment-id: o014.li2.chapter4.morphisms-extensions.g001
	\begin{equation}\label{eqn:Hom-D-injective-aux}
		\Hom_{\cate{D}(\mathcal{A})}(T, X) \simeq \varinjlim_{\substack{\alpha: S \to T \\ \text{quasi-isomorphism in }\cate{K}(\mathcal{A})}} \Hom_{\cate{K}(\mathcal{A})}(S, X), \quad T \in \Obj(\cate{K}(\mathcal{A})).
	\end{equation}

	Suppose $X \in \Obj(\cate{C}^+(\mathcal{A}))$ consists of injective
	objects. For a quasi-isomorphism
	$\alpha \in \Hom_{\cate{K}(\mathcal{A})}(S, T)$,
	Theorem~\ref{prop:resolution-homotopy} shows that
	$\alpha^*: \Hom_{\cate{K}(\mathcal{A})}(T, X) \rightiso
	\Hom_{\cate{K}(\mathcal{A})}(S, X)$. Hence the $\varinjlim$ in
	\eqref{eqn:Hom-D-injective-aux} is equal to
	$\Hom_{\cate{K}(\mathcal{A})}(T, X)$.

	If $X \in \Obj(\cate{C}(\mathcal{A}))$ is K-injective, use
	Theorem~\ref{prop:K-resolution-homotopy} instead.
\end{proof}

% segment-id: o014.li2.chapter4.morphisms-extensions.p002
One immediate application of the preceding proposition is to control the size
of a derived category. Here we again use the chosen Grothendieck universe
$\mathcal{U}$, together with the terminology of $\mathcal{U}$-categories
(that is, “categories” in the standard convention of this book) and small
$\mathcal{U}$-categories.

% segment-id: o014.li2.chapter4.morphisms-extensions.co001
\begin{corollary}\label{prop:D-size}
	Use the notation above. Let $\mathcal{A}$ be a $\mathcal{U}$-category.
% segment-id: o014.li2.chapter4.morphisms-extensions.l001
	\begin{enumerate}[(i)]
% segment-id: o014.li2.chapter4.morphisms-extensions.i001
		\item If $\mathcal{A}$ has enough injective objects (respectively
		projective objects), then $\cate{D}^+(\mathcal{A})$ (respectively
		$\cate{D}^-(\mathcal{A})$) is also a $\mathcal{U}$-category.
% segment-id: o014.li2.chapter4.morphisms-extensions.i002
		\item If $\mathcal{A}$ has enough K-injective or K-projective
		complexes (Definition~\ref{def:K-resolutions}), then
		$\cate{D}(\mathcal{A})$ is also a $\mathcal{U}$-category.
% segment-id: o014.li2.chapter4.morphisms-extensions.i003
		\item If $\mathcal{A}$ is a small $\mathcal{U}$-category, then so is
		$\cate{D}(\mathcal{A})$.
	\end{enumerate}
\end{corollary}

% segment-id: o014.li2.chapter4.morphisms-extensions.pf002
\begin{proof}
	It suffices to verify that $\Hom_{\cate{D}(\mathcal{A})}(X,Y)$ is a small
	$\mathcal{U}$-set. Notice that the $\Hom_{\cate{K}(\mathcal{A})}$ version
	is easy to handle, since it is straightforward to prove that
	$\cate{C}(\mathcal{A})$ is a $\mathcal{U}$-category.

	Suppose $\mathcal{A}$ has enough injective objects and
	$X \in \Obj(\cate{D}^+(\mathcal{A}))$. Without loss of generality, we
	may assume that $X \in \Obj(\cate{C}^+(\mathcal{A}))$ consists of
	injective objects; (i) then follows from
	Proposition~\ref{prop:Hom-D-injective}. If $X$ has a K-injective resolution
	(respectively $Y$ has a K-projective resolution), without loss of
	generality we may assume that $X$ is K-injective (respectively $Y$ is
	K-projective); the same argument gives (ii).

	For (iii), observe that $\cate{K}(\mathcal{A})$ is then a small
	$\mathcal{U}$-category. Apply this observation to
	Remark~\ref{rem:localization-size}.
\end{proof}

% segment-id: o014.li2.chapter4.morphisms-extensions.x001
\index{size issues}

% segment-id: o014.li2.chapter4.morphisms-extensions.pr002
\begin{proposition}[Orthogonality]\label{prop:derived-cat-orthogonality}
	Let $a,b,n\in\Z$, with $a<b$.
% segment-id: o014.li2.chapter4.morphisms-extensions.l002
	\begin{enumerate}[(i)]
% segment-id: o014.li2.chapter4.morphisms-extensions.i004
		\item If $X \in \Obj(\cate{D}^{\leq a}(\mathcal{A}))$ and
		$Y \in \Obj(\cate{D}^{\geq b}(\mathcal{A}))$, then
		$\Hom_{\cate{D}(\mathcal{A})}(X,Y)=0$.
% segment-id: o014.li2.chapter4.morphisms-extensions.i005
		\item If $X \in \Obj(\cate{D}^{\leq n}(\mathcal{A}))$ and
		$Y \in \Obj(\cate{D}^{\geq n}(\mathcal{A}))$, then $\Hm^n$
		induces a canonical isomorphism
% segment-id: o014.li2.chapter4.morphisms-extensions.q001
		\[ \Hom_{\cate{D}(\mathcal{A})}(X, Y) \rightiso \Hom_{\mathcal{A}}\left( \Hm^n(X), \Hm^n(Y) \right). \]
	\end{enumerate}
\end{proposition}

% segment-id: o014.li2.chapter4.morphisms-extensions.pf003
\begin{proof}
	Let $a\leq b$. Represent a morphism
	$f\in\Hom_{\cate{D}(\mathcal{A})}(X,Y)$ by a diagram
	$X\xleftarrow{s}Z\xrightarrow{g}Y$ in $\cate{K}(\mathcal{A})$, where
	$s$ is a quasi-isomorphism. Replacing $f$ by $g$, we may reduce to the
	case in which $f$ comes from $\cate{C}(\mathcal{A})$. Next replace $f$
	by the composite
	$\tau^{\leq a}X\to X\xrightarrow{f}Y\to\tau^{\geq b}Y$ in
	$\cate{C}(\mathcal{A})$. The two morphisms at the ends of this composite
	are quasi-isomorphisms. We may therefore reduce further to the case
	$X\in\Obj(\cate{C}^{\leq a}(\mathcal{A}))$ and
	$Y\in\Obj(\cate{C}^{\geq b}(\mathcal{A}))$.

	If $a<b$, every morphism $X\to Y$ in $\cate{C}(\mathcal{A})$ must be
	zero. This proves (i).

	Now take $X$ and $Y$ as above, but set $a=n=b$ to treat (ii). Consider
	the canonical morphisms
% segment-id: o014.li2.chapter4.morphisms-extensions.e001
	\begin{equation}\label{eqn:derived-cat-orthogonality-aux}
		\Hom_{\mathcal{A}}\left( \Hm^n(X), \Hm^n(Y) \right) \xleftarrow{\Hm^n} \Hom_{\cate{C}(\mathcal{A})}(X, Y) \to \Hom_{\cate{K}(\mathcal{A})}(X, Y) \to \Hom_{\cate{D}(\mathcal{A})}(X, Y).
	\end{equation}
	\begin{itemize}
		\item The first morphism $\Hm^n$ is an isomorphism: because
		$X\in\Obj(\cate{C}^{\leq n}(\mathcal{A}))$ and
		$Y\in\Obj(\cate{C}^{\geq n}(\mathcal{A}))$, any
		$f\in\Hom_{\cate{C}(\mathcal{A})}(X,Y)$ can be nonzero only in
		degree $n$, and the only conditions on $f^n$ are
% segment-id: o014.li2.chapter4.morphisms-extensions.q002
		\[ f^n\left( \Image(d_X^{n-1}) \right) = 0 , \quad f\left(X^n\right) \subset \Ker\left( d_Y^n \right). \]
		But $X^n/\Image(d_X^{n-1})=\Hm^n(X)$ and
		$\Ker(d_Y^n)=\Hm^n(Y)$.
		\item The second morphism is also an isomorphism, since the conditions
		on $X,Y$ imply $\Hom^{-1}(X,Y)=\{0\}$.
	\end{itemize}

	Thus, to prove (ii), it remains to show that the last morphism in
	\eqref{eqn:derived-cat-orthogonality-aux} is an isomorphism. Applying
	Lemma~\ref{prop:MXY-equiv} gives
% segment-id: o014.li2.chapter4.morphisms-extensions.q003
	\[ \Hom_{\cate{D}(\mathcal{A})}(X, Y) \simeq \varinjlim_{[Y \rightarrowtail Z] \in \Obj(\mathrm{qis}_{Y/}) } \Hom_{\cate{K}(\mathcal{A})}(X, Z). \]
	If $Y\to Z$ is a quasi-isomorphism, then both morphisms in
	$Y\to Z\to\tau^{\geq n}Z$ are quasi-isomorphisms in
	$\cate{K}(\mathcal{A})$. This shows that the quasi-isomorphisms $Y\to Z$
	with $Z$ coming from $\cate{C}^{\geq n}(\mathcal{A})$ form a cofinal full
	subcategory of the filtered category $\mathrm{qis}_{Y/}$
	(Definition~\ref{def:cofinal} and
	Proposition~\ref{prop:cofinal-filtered}). By
	Proposition~\ref{prop:cofinal-lim}, we may henceforth restrict the
	$\varinjlim$ to this subcategory.

	We already know that the first two morphisms in
	\eqref{eqn:derived-cat-orthogonality-aux} are canonical isomorphisms. Take
	$Y\to Z$ as above. This gives the commutative diagram
% segment-id: o014.li2.chapter4.morphisms-extensions.g002
	\[\begin{tikzcd}
		\Hom_{\mathcal{A}}(\Hm^n(X), \Hm^n(Y)) \arrow[d, "\sim" sloped] & \Hom_{\cate{K}(\mathcal{A})}(X, Y) \arrow[l, "\sim", "{\Hm^n}"'] \arrow[d] \\
		\Hom_{\mathcal{A}}(\Hm^n(X), \Hm^n(Z)) & \Hom_{\cate{K}(\mathcal{A})}(X, Z) \arrow[l, "\sim"', "{\Hm^n}"]
	\end{tikzcd}\]
	which shows that
	$\Hom_{\cate{K}(\mathcal{A})}(X,Y)\rightiso
	\Hom_{\cate{K}(\mathcal{A})}(X,Z)$. Hence the preceding $\varinjlim$ is
	constant with value $\Hom_{\cate{K}(\mathcal{A})}(X,Y)$. The last
	morphism in \eqref{eqn:derived-cat-orthogonality-aux} is therefore an
	isomorphism.
\end{proof}

% segment-id: o014.li2.chapter4.morphisms-extensions.th001
\begin{theorem}\label{prop:derived-cat-embedding}
	Regarding the objects of $\mathcal{A}$ as complexes concentrated in
	degree zero gives an equivalence of additive categories
	$\mathcal{A}\to\cate{D}^{\leq0}(\mathcal{A})\cap
	\cate{D}^{\geq0}(\mathcal{A})$, with quasi-inverse given by $\Hm^0$.
\end{theorem}

% segment-id: o014.li2.chapter4.morphisms-extensions.pf004
\begin{proof}
	It is clear that the objects of $\mathcal{A}$ give objects of
	$\cate{D}^{\leq0}(\mathcal{A})\cap\cate{D}^{\geq0}(\mathcal{A})$.
	Thus we obtain an additive functor
	$\Phi:\mathcal{A}\to\cate{D}^{\leq0}(\mathcal{A})\cap
	\cate{D}^{\geq0}(\mathcal{A})$. Proposition
	\ref{prop:derived-cat-orthogonality} (ii) implies that $\Phi$ is fully
	faithful.

	We next show that $\Phi$ is essentially surjective. Let $X$ be an object
	of $\cate{D}^{\leq0}(\mathcal{A})\cap\cate{D}^{\geq0}(\mathcal{A})$.
	Applying the truncation functors to the complex $X$ in
	$\cate{C}(\mathcal{A})$, we see that every morphism in
	$X\to\tau^{\geq0}X\leftarrow\tau^{\leq0}\tau^{\geq0}X$ is a
	quasi-isomorphism, while Proposition~\ref{prop:truncate-twice} gives
	$\tau^{\leq0}\tau^{\geq0}X\simeq\Hm^0(X)$. Hence $\Phi$ is an
	equivalence and $\Hm^0$ is its quasi-inverse.
\end{proof}

% segment-id: o014.li2.chapter4.morphisms-extensions.p003
More generally, for every $n\in\Z$, the functor $S\mapsto S[-n]$ gives an
additive equivalence from $\mathcal{A}$ to
$\cate{D}^{\geq n}(\mathcal{A})\cap\cate{D}^{\leq n}(\mathcal{A})$, with
$\Hm^n$ as quasi-inverse. This is simply the shifted version of the result
above.

% segment-id: o014.li2.chapter4.morphisms-extensions.p004
Henceforth, without further comment, we identify $\mathcal{A}$ with the full
additive subcategory
$\cate{D}^{\geq0}(\mathcal{A})\cap\cate{D}^{\leq0}(\mathcal{A})$ of
$\cate{D}^{\bdd}(\mathcal{A})$.

% segment-id: o014.li2.chapter4.morphisms-extensions.d001
\begin{definition}\label{def:functor-amplitude}
	\index{amplitude}
	Let $\star\in\{+,-,\bdd,\hspace{0.8em}\}$ (where
	$\star=\hspace{0.8em}$ means that there is no superscript), and let
	$\mathcal{A}$ and $\mathcal{A}'$ be Abelian categories. If a
	triangulated functor
	$R:\cate{D}^\star(\mathcal{A})\to\cate{D}^\star(\mathcal{A}')$
	restricts to
	$\mathcal{A}\to\cate{D}^{[a,b]}(\mathcal{A}')$, where
	$-\infty\leq a\leq b\leq+\infty$, then $R$ is said to have
	\emph{amplitude} contained in $[a,b]$.
\end{definition}

% segment-id: o014.li2.chapter4.morphisms-extensions.p005
Here is a simple observation: if
$X'\to X\to X''\xrightarrow{+1}$ is a distinguished triangle in
$\cate{D}(\mathcal{A})$ and
$X',X''\in\Obj(\cate{D}^{[a,b]}(\mathcal{A}))$, then
$X\in\Obj(\cate{D}^{[a,b]}(\mathcal{A}))$. This follows immediately from
the long exact sequence for $\Hm^n$.

% segment-id: o014.li2.chapter4.morphisms-extensions.le001
\begin{lemma}\label{prop:amplitude-translate}
	If the amplitude of $R$ is contained in $[a,b]$, then for every
	$-\infty<c\leq d<+\infty$ one has
% segment-id: o014.li2.chapter4.morphisms-extensions.q004
	\[ R\;\text{restricts to}\;\cate{D}^{[c,d]}(\mathcal{A})\to\cate{D}^{[c+a,d+b]}(\mathcal{A}'). \]
\end{lemma}

% segment-id: o014.li2.chapter4.morphisms-extensions.pf005
\begin{proof}
	The case $c=d$ is immediate. If $c<d$, then for every
	$X\in\Obj(\cate{D}^{[c,d]}(\mathcal{A}))$, take the distinguished
	triangle from Proposition~\ref{prop:derived-truncation-adjunction}
% segment-id: o014.li2.chapter4.morphisms-extensions.q005
	\[ \tau^{\leq d-1}X\to X\to\tau^{\geq d}X\xrightarrow{+1},\quad \tau^{\geq d}X\simeq\Hm^d(X)[-d] \]
	and apply the preceding observation to its image under $R$ to carry out
	the induction.
\end{proof}

% segment-id: o014.li2.chapter4.morphisms-extensions.p006
Consequently, if $R$ has finite amplitude, then it restricts to
$\cate{D}^{\bdd}(\mathcal{A})\to\cate{D}^{\bdd}(\mathcal{A}')$.

% segment-id: o014.li2.chapter4.morphisms-extensions.p007
We next give a related result explaining how a triangulated functor is
determined by its values on $\mathcal{A}$.

% segment-id: o014.li2.chapter4.morphisms-extensions.pr003
\begin{proposition}[Way-out Lemma {\cite[Chapter I, \S 7]{Har66}}]\label{prop:wayout}
	\index{Way-out Lemma}
	Let $\mathcal{T}$ be a weak Serre subcategory of $\mathcal{A}$ and let
	$\star\in\{\hspace{0.8em},+,\bdd\}$. Consider triangulated functors
	$F,G:\cate{D}^{\star}_{\mathcal{T}}(\mathcal{A})\to
	\cate{D}(\mathcal{A}')$ and a morphism $\eta:F\to G$ between them such
	that $\eta_X:FX\to GX$ is an isomorphism for every
	$X\in\Obj(\mathcal{T})$. Then $\eta$ is also an isomorphism if one of the
	following conditions holds:
% segment-id: o014.li2.chapter4.morphisms-extensions.l003
	\begin{enumerate}[(i)]
% segment-id: o014.li2.chapter4.morphisms-extensions.i006
		\item $\star=\bdd$;
% segment-id: o014.li2.chapter4.morphisms-extensions.i007
		\item $\star=+$, and there exists $k\in\Z$ such that $F,G$ restrict
		to $\cate{D}^{\geq0}_{\mathcal{T}}(\mathcal{A})\to
		\cate{D}^{\geq k}(\mathcal{A}')$;
% segment-id: o014.li2.chapter4.morphisms-extensions.i008
		\item $\star=\hspace{0.8em}$, and there exist $k,\ell\in\Z$ such that
		$F,G$ respectively restrict to
		$\cate{D}^{\geq0}_{\mathcal{T}}(\mathcal{A})\to
		\cate{D}^{\geq k}(\mathcal{A}')$ and
		$\cate{D}^{\leq0}_{\mathcal{T}}(\mathcal{A})\to
		\cate{D}^{\leq\ell}(\mathcal{A}')$.
	\end{enumerate}

	In addition, let $\mathfrak{I}$ be a subset of $\Obj(\mathcal{T})$ such
	that for every $X\in\Obj(\mathcal{T})$ there exist
	$I\in\mathfrak{I}$ and a monomorphism $X\hookrightarrow I$. If the
	condition on $\eta_X$ is replaced by the condition that $\eta_I$ be an
	isomorphism for every $I\in\mathfrak{I}$, then $\eta$ is still an
	isomorphism under assumption (i) or (ii).

	Of course, (ii) also has a dual version for $\star=-$, which we do not
	write out.
\end{proposition}

% segment-id: o014.li2.chapter4.morphisms-extensions.pf006
\begin{proof}
	For (i), simply repeat the argument of
	Lemma~\ref{prop:amplitude-translate}; use the distinguished triangle
	$\tau^{\leq d}X\to X\to\tau^{\geq d+1}X\xrightarrow{+1}$ and
	Proposition~\ref{prop:dt-3-2} inductively to reduce to the case
	$X\simeq\Hm^{d+1}(X)[-d-1]$.

	For (ii), the goal is to prove that $\Hm^j(FX)\to\Hm^j(GX)$ is an
	isomorphism for every
	$X\in\Obj(\cate{D}^+_{\mathcal{T}}(\mathcal{A}))$ and $j\in\Z$.
	Consider the distinguished triangle above once more, but take $d\gg0$ so
	that $F\tau^{\geq d+1}X$ and $G\tau^{\geq d+1}X$ both lie in
	$\cate{D}^{\geq j+1}(\mathcal{A}')$. We obtain a commutative diagram with
	exact rows
% segment-id: o014.li2.chapter4.morphisms-extensions.g003
	\[\begin{tikzcd}
		0 \arrow[r] \arrow[d] & \Hm^j\left(F\tau^{\leq d}X\right) \arrow[r] \arrow[d] & \Hm^j(FX) \arrow[d] \arrow[r] & 0 \arrow[d] \\
		0 \arrow[r] & \Hm^j\left(G\tau^{\leq d}X\right) \arrow[r] & \Hm^j(GX) \arrow[r] & 0
	\end{tikzcd}\]
	But $\tau^{\leq d}X\in
	\Obj(\cate{D}^{\bdd}_{\mathcal{T}}(\mathcal{A}))$, so this case
	follows from (i).

	For (iii), choose any $d\in\Z$. Using
	$\tau^{\leq d}X\to X\to\tau^{\geq d+1}X\xrightarrow{+1}$ and
	Proposition~\ref{prop:dt-3-2}, this case follows from (ii) and its dual.

	Now consider the final assertion. Since (ii) has been proved, it suffices
	to show that $\eta_X$ is an isomorphism for every
	$X\in\Obj(\mathcal{T})$. Take an exact sequence
	$0\to X\to I^0\to I^1\to\cdots$ with every $I^n\in\mathfrak{I}$.
	Write this sequence as a quasi-isomorphism of complexes $X\to I$; it
	remains to prove that $\eta_I$ is an isomorphism. To do this, define the
	brutal truncation $\sigma^{\leq d}I$ of the complex, which agrees with
	$I$ in degrees $\leq d$ and has zero terms in all other degrees; define
	$\sigma^{\geq d+1}I$ similarly. We can now follow the arguments for (i)
	and (ii), using instead the canonical short exact sequence in
	$\cate{C}(\mathcal{A})$
% segment-id: o014.li2.chapter4.morphisms-extensions.q006
	\[ 0\to\sigma^{\geq d+1}I\to I\to\sigma^{\leq d}I\to0 \]
	and its corresponding distinguished triangle. The details are left to the
	reader.
\end{proof}

% segment-id: o014.li2.chapter4.morphisms-extensions.p008
Recall Theorem~\ref{prop:derived-cat-embedding}. Its key ingredient is the
full faithfulness
$\Hom_{\mathcal{A}}(X,Y)\simeq\Hom_{\cate{D}(\mathcal{A})}(X,Y)$.
What information do we obtain when the degrees of $X$ and $Y$ are shifted
relative to one another? This is our next topic.

% segment-id: o014.li2.chapter4.morphisms-extensions.d002
\begin{definition}[$\Ext$ functors: General case]\label{def:Ext-general}
	\index[sym1]{Extn}
	\index{$\Ext$ functor}
	For $X,Y\in\Obj(\mathcal{A})$ and $n\in\Z$, define
	$\Ext^n_{\mathcal{A}}(X,Y):=
	\Hom_{\cate{D}(\mathcal{A})}\left(X,Y[n]\right)$. The functor
	$\Ext^n_{\mathcal{A}}:\mathcal{A}^{\opp}\times\mathcal{A}\to
	\cate{Ab}$ is additive in each variable.
\end{definition}

% segment-id: o014.li2.chapter4.morphisms-extensions.p009
The notation $\Ext^n_{\mathcal{A}}(X,Y)$ already appeared in
Definition--Proposition~\ref{def:Ext-classical}. The discussion of $\RHom$
will reconcile the two definitions; see
Corollary~\ref{prop:Ext-reconcilation}. Here we do not assume that
$\mathcal{A}$ has enough injective or projective objects.

% segment-id: o014.li2.chapter4.morphisms-extensions.p010
As for $\Hom$, functoriality of $\Ext^n$ in its two variables is still called
pullback and pushout, respectively. When no confusion can arise,
$\Ext^n_{\mathcal{A}}$ is also abbreviated to $\Ext^n$. We now show that its
properties resemble those of the $\Ext^n$ defined in \S\ref{sec:Ext-Tor}.

% segment-id: o014.li2.chapter4.morphisms-extensions.pr004
\begin{proposition}\label{prop:Ext-generalities}
	Below, $X,Y$ are arbitrary objects of $\mathcal{A}$. The family of
	bifunctors $(\Ext^n)_{n\in\Z}$ has the following properties.
% segment-id: o014.li2.chapter4.morphisms-extensions.l004
	\begin{enumerate}[(i)]
% segment-id: o014.li2.chapter4.morphisms-extensions.i009
		\item If $n<0$, then $\Ext^n(X,Y)=0$.
% segment-id: o014.li2.chapter4.morphisms-extensions.i010
		\item There is a canonical isomorphism
		$\Ext^0(X,Y)\simeq\Hom_{\mathcal{A}}(X,Y)$.
% segment-id: o014.li2.chapter4.morphisms-extensions.i011
		\item Short exact sequences
		$0\to X'\to X\to X''\to0$ and
		$0\to Y'\to Y\to Y''\to0$ in $\mathcal{A}$ induce the respective
		long exact sequences
% segment-id: o014.li2.chapter4.morphisms-extensions.q007
		\begin{gather*}
			\cdots \to \Ext^{n-1}(X', Y) \xrightarrow{\delta^{n-1}} \Ext^n(X'', Y) \to \Ext^n(X, Y) \to \Ext^n(X', Y) \to \cdots , \\
			\cdots \to \Ext^{n-1}(X, Y'') \xrightarrow{\delta^{n-1}} \Ext^n(X, Y') \to \Ext^n(X, Y) \to \Ext^n(X, Y'') \to \cdots ,
		\end{gather*}
		with every connecting morphism $\delta^n$ functorial in the short
		exact sequence.
% segment-id: o014.li2.chapter4.morphisms-extensions.i012
		\item If $\Ext^1(X,\cdot)=0$ (respectively
		$\Ext^1(\cdot,Y)=0$), then $X$ (respectively $Y$) is a projective
		(respectively injective) object of $\mathcal{A}$.
	\end{enumerate}
\end{proposition}

% segment-id: o014.li2.chapter4.morphisms-extensions.pf007
\begin{proof}
	Assertions (i) and (ii) follow immediately from
	Proposition~\ref{prop:derived-cat-orthogonality}.

	For (iii), first use Proposition~\ref{prop:ses-vs-triangle} to turn the
	short exact sequences canonically into distinguished triangles
	$X'\to X\to X''\xrightarrow{+1}$ and
	$Y'\to Y\to Y''\xrightarrow{+1}$ in $\cate{D}(\mathcal{A})$. Then use
	the fact that $\Hom$ is a cohomological functor in each variable
	(Proposition~\ref{prop:cohomological-Hom-functor}).

	Finally, (iv) follows by applying (i) and (ii) to the long exact sequences
	in (iii).
\end{proof}

% segment-id: o014.li2.chapter4.morphisms-extensions.p011
Since $\Ext^n$ is given by $\Hom$ in $\cate{D}(\mathcal{A})$, for
$f\in\Ext^n(X,Y)$ and $g\in\Ext^n(X',Y')$ there is an evident direct sum
$f\oplus g\in\Ext^n(X\oplus X',Y\oplus Y')$. On the other hand, $\Ext$
has a composition operation,
% segment-id: o014.li2.chapter4.morphisms-extensions.g004
\begin{equation}\label{eqn:Ext-composition}\begin{tikzcd}[row sep=tiny]
	\Ext^m(Y, Z) \times \Ext^n(X, Y) \arrow[equal, d] & \\
	\Hom_{\cate{D}(\mathcal{A})}(Y, Z[m]) \times \Hom_{\cate{D}(\mathcal{A})}(X, Y[n]) \arrow[r] & \Ext^{m+n}(X, Z) \\
	(f, g) \arrow[r, mapsto] & fg := f[n] \circ g ,
\end{tikzcd}\end{equation}
where $m,n\in\Z$.

% segment-id: o014.li2.chapter4.morphisms-extensions.p012
It is easy to see that this operation is associative and bilinear.
Consequently, the Abelian group
% segment-id: o014.li2.chapter4.morphisms-extensions.q008
\[ \Ext(X):=\bigoplus_{n\in\Z}\Ext^n(X,X) \]
has a natural ring structure, with
$\identity_A\in\End_{\mathcal{A}}(A)$ as multiplicative identity, while
$\Ext(X,Y):=\bigoplus_{n\in\Z}\Ext^n(X,Y)$ becomes an
$(\Ext(Y),\Ext(X))$-bimodule.

% segment-id: o014.li2.chapter4.morphisms-extensions.p013
More generally, if $\mathcal{A}$ is $\Bbbk$-linear, where $\Bbbk$ is a
commutative ring, then the functors $\Ext^n$ take values in
$\Bbbk\dcate{Mod}$, while $\Ext(X)$ becomes a $\Bbbk$-algebra. The
definition of multiplication in \eqref{eqn:Ext-composition} shows that
$\Ext(X)$ is in fact a graded $\Bbbk$-algebra; see
\cite[Definition~7.4.1]{Li1}. Such constructions are extremely useful in
representation theory.

% segment-id: o014.li2.chapter4.morphisms-extensions.d003
\begin{definition}\index{$\Ext$-algebra}\label{def:Ext-algebra}
	Let $X\in\Obj(\mathcal{A})$. The graded algebra $\Ext(X)$ defined above
	is called the \emph{$\Ext$-algebra} of the object $X$.
\end{definition}

% segment-id: o014.li2.chapter4.morphisms-extensions.p014
The notation $\Ext$ abbreviates “extension.” To explain the concrete
connection, we first define what an extension means.

% segment-id: o014.li2.chapter4.morphisms-extensions.d004
\begin{definition}\label{def:Yoneda-Ext}
	\index{extension}
	Let $X$ and $Y$ be objects of $\mathcal{A}$ and let
	$n\in\Z_{\geq1}$. An exact sequence in $\mathcal{A}$ of the form
% segment-id: o014.li2.chapter4.morphisms-extensions.q009
	\[ \mathcal{E}:\quad 0\to Y\to E^1\to\cdots\to E^n\to X\to0 \]
	is called an \emph{$n$-extension} of $X$ by $Y$. For fixed $X,Y$,
	consider the following binary relation on the set of all $n$-extensions:
% segment-id: o014.li2.chapter4.morphisms-extensions.g005
	\begin{multline*} \mathcal{E} \sim_0 \mathcal{E}' \iff \text{there is a commutative diagram} \\
		\begin{tikzcd}[ampersand replacement=\&]
			0 \arrow[r] \& Y \arrow[r] \arrow[equal, d] \& E^1 \arrow[r] \arrow[d] \& \cdots \arrow[r] \& E^n \arrow[d] \arrow[r] \& X \arrow[equal, d] \arrow[r] \& 0 \\
			0 \arrow[r] \& Y \arrow[r] \& (E')^1 \arrow[r] \& \cdots \arrow[r] \& (E')^n \arrow[r] \& X \arrow[r] \& 0
		\end{tikzcd}\end{multline*}
	This binary relation generates an equivalence relation $\sim$.\footnote{In
	other words, $\mathcal{E}_1\sim\mathcal{E}_2$ if and only if they can be
	connected by a chain of commutative diagrams of this kind.} The
	equivalence classes of $n$-extensions form the set
	$\Ext^{n,\text{Yoneda}}(X,Y)$.
\end{definition}

% segment-id: o014.li2.chapter4.morphisms-extensions.p015
Proposition~\ref{prop:5-lemma} implies that equivalence of $1$-extensions is
the same as isomorphism of short exact sequences
% segment-id: o014.li2.chapter4.morphisms-extensions.g006
\[\begin{tikzcd}
	0 \arrow[r] & Y \arrow[r] \arrow[equal, d] & E \arrow[r] \arrow[d, "\sim" sloped] & X \arrow[equal, d] \arrow[r] & 0 \\
	0 \arrow[r] & Y \arrow[r] & E' \arrow[r] & X \arrow[r] & 0 .
\end{tikzcd}\]

% segment-id: o014.li2.chapter4.morphisms-extensions.cv001
\begin{convention}
	Henceforth a $1$-extension is simply called an extension, and equivalence
	between such extensions is called isomorphism. A split short exact
	sequence (Proposition~\ref{prop:split-ses}) is also called a split
	extension; all such extensions are isomorphic.
\end{convention}

% segment-id: o014.li2.chapter4.morphisms-extensions.p016
The following operations are available on $n$-extensions.

% segment-id: o014.li2.chapter4.morphisms-extensions.l005
\begin{description}
% segment-id: o014.li2.chapter4.morphisms-extensions.i013
	\item[Composition] For
	$[\mathcal{E}_1]\in\Ext^{n,\text{Yoneda}}(Y,Z)$ and
	$[\mathcal{E}_2]\in\Ext^{m,\text{Yoneda}}(X,Y)$, define
	$[\mathcal{E}_1]\circ[\mathcal{E}_2]\in
	\Ext^{n+m,\text{Yoneda}}(X,Z)$ by splicing the two exact sequences
	end-to-end at $Y$. This operation is also called the \emph{Yoneda
	product}.
	\index{extension!Yoneda product}
% segment-id: o014.li2.chapter4.morphisms-extensions.i014
	\item[Pullback] For a morphism $f:X'\to X$, the corresponding map
	$f^*:\Ext^{n,\text{Yoneda}}(X,Y)\to
	\Ext^{n,\text{Yoneda}}(X',Y)$ is obtained as follows. Given
	$[\mathcal{E}]\in\Ext^{n,\text{Yoneda}}(X,Y)$, consider the commutative
	diagram
% segment-id: o014.li2.chapter4.morphisms-extensions.g007
	\[\begin{tikzcd}[every cell/.append style={font=\small}]
		0 \arrow[r] & Y \arrow[r] \arrow[equal, d] & E^1 \arrow[r] \arrow[equal, d] & \cdots \arrow[r] & E^{n-1}  \arrow[equal, d] \arrow[r] & E^n \dtimes{X} X' \arrow[d] \arrow[r] \arrow[phantom, rd, "\Box" description] & X' \arrow[d, "f"] \arrow[r] & 0 \\
		0 \arrow[r] & Y \arrow[r] & E^1 \arrow[r] & \cdots \arrow[r] & E^{n-1} \arrow[r] & E^n \arrow[r] & X \arrow[r] & 0
	\end{tikzcd}\]
	The square marked $\Box$ is a pullback square, while
	$E^{n-1}\to E^n\dtimes{X}X'$ is determined by $E^{n-1}\to E^n$ and
	$E^{n-1}\xrightarrow{0}X'$. Pullback preserves kernels, and
	$E^n\dtimes{X}X'\to X'$ is still an epimorphism
	(Proposition~\ref{prop:Abel-cat-pull-push}); from this one checks that the
	first row remains exact. Its equivalence class is $f^*[\mathcal{E}]$.

% segment-id: o014.li2.chapter4.morphisms-extensions.i015
	\item[Pushout] For a morphism $g:Y\to Y'$, the corresponding map
	$g_*:\Ext^{n,\text{Yoneda}}(X,Y)\to
	\Ext^{n,\text{Yoneda}}(X,Y')$ is obtained as follows. Given
	$[\mathcal{E}]\in\Ext^{n,\text{Yoneda}}(X,Y)$, pushout gives the
	commutative diagram
% segment-id: o014.li2.chapter4.morphisms-extensions.g008
	\[\begin{tikzcd}[every cell/.append style={font=\small}]
		0 \arrow[r] & Y \arrow[r] \arrow[d, "g"'] \arrow[phantom, rd, "\boxplus" description] & E^1 \arrow[r] \arrow[d] & E^2 \arrow[r] \arrow[equal, d] & \cdots \arrow[r] & E^n \arrow[equal, d] \arrow[r] & X \arrow[equal, d] \arrow[r] & 0 \\
		0 \arrow[r] & Y' \arrow[r] & E^1 \dsqcup{Y} Y' \arrow[r] & E^2 \arrow[r] & \cdots \arrow[r] & E^n \arrow[r] & X \arrow[r] & 0
	\end{tikzcd}\]
	Here $E^1\dsqcup{Y}Y'\to E^2$ is determined by $E^1\to E^2$ and
	$Y'\xrightarrow{0}E^2$. Dually, one checks that the second row remains
	exact. Its equivalence class is $g_*[\mathcal{E}]$.

% segment-id: o014.li2.chapter4.morphisms-extensions.i016
	\item[Direct sum] Take the termwise direct sum of two extensions,
	$\mathcal{E}\oplus\mathcal{E}'$, and then take its equivalence class.
	This gives a map
% segment-id: o014.li2.chapter4.morphisms-extensions.g009
	\[\begin{tikzcd}
		\oplus: \Ext^{n, \text{Yoneda}}(X, Y) \times \Ext^{n, \text{Yoneda}}(X', Y') \arrow[r] & \Ext^{n, \text{Yoneda}}(X \oplus X', Y \oplus Y') .
	\end{tikzcd}\]

% segment-id: o014.li2.chapter4.morphisms-extensions.i017
	\item[Baer sum] This is the following map:
	\index{extension!Baer sum}
	\index[sym1]{plusdot@$\dot{+}$}
% segment-id: o014.li2.chapter4.morphisms-extensions.g010
	\[\begin{tikzcd}
		\dot{+}: \Ext^{n, \text{Yoneda}}(X, Y) \times \Ext^{n, \text{Yoneda}}(X, Y) \arrow[r] & \Ext^{n, \text{Yoneda}}(X, Y).
	\end{tikzcd}\]
	It is defined by pulling $[\mathcal{E}_1]\oplus[\mathcal{E}_2]$ back
	along $X\hookrightarrow X\times X$, then pushing it out along
	$Y\oplus Y\twoheadrightarrow Y$ (these morphisms come from
	Convention~\ref{con:biproduct}). The result is
	$[\mathcal{E}_1]\dot{+}[\mathcal{E}_2]\in
	\Ext^{n,\text{Yoneda}}(X,Y)$.
\end{description}

% segment-id: o014.li2.chapter4.morphisms-extensions.p017
One checks that pullback and pushout commute:
$g_*f^*=f^*g_*$, while the Baer sum $\dot{+}$ is commutative. These
properties also follow from those of $\Ext^n$ by the following theorem.

% segment-id: o014.li2.chapter4.morphisms-extensions.th002
\begin{theorem}[Nobuo Yoneda]\label{prop:Yoneda-Ext}
	For all $X,Y\in\Obj(\mathcal{A})$ and $n\in\Z_{\geq1}$, there is a
	canonical bijection
% segment-id: o014.li2.chapter4.morphisms-extensions.q010
	\[ \Ext^{n,\text{Yoneda}}(X,Y) \xrightarrow{1:1} \Ext^n(X,Y), \]
	which sends the equivalence class of an $n$-extension
	$0\to Y\to E^1\to\cdots\to E^n\to X\to0$ to
	$as^{-1}\in\Hom_{\cate{D}(\mathcal{A})}(X,Y[n])$, where
% segment-id: o014.li2.chapter4.morphisms-extensions.g011
	\[\begin{tikzcd}
		X \arrow[phantom, r, "" {name=A, yshift=2ex}] & {} & & & X \\
		Z \arrow[u, "{s: \text{quasi-isomorphism}}"] \arrow[d, "a"'] & Y \arrow[r] \arrow[d, "\identity"'] & E^1 \arrow[r] & \cdots \arrow[r] & E^n \arrow[u] \\
		Y[n] \arrow[phantom, r, "" {name=B, yshift=-2ex}] & Y & & &
		\arrow[dash, from=A, to=B]
	\end{tikzcd}\]
	and all omitted terms are zero. Under this bijection, composition,
	direct sum, pullback, and pushout of $n$-extensions correspond to the
	respective operations on $\Ext^n$, while the Baer sum $\dot{+}$
	corresponds to the additive group law on $\Ext^n(X,Y)$.
\end{theorem}

% segment-id: o014.li2.chapter4.morphisms-extensions.pf008
\begin{proof}
	An equivalence of extensions $\mathcal{E}\sim\mathcal{E}'$ is reflected
	by a quasi-isomorphism between the corresponding complexes $Z$ and
	$Z'$. Thus the displayed map is well defined. It remains to prove that
	it is surjective and injective. We begin with surjectivity.

	Every $f\in\Hom_{\cate{D}(\mathcal{A})}(X,Y[n])$ is represented by a
	diagram $X\xleftarrow{s}Z\xrightarrow{a}Y[n]$ in
	$\cate{C}(\mathcal{A})$, where $s$ is a quasi-isomorphism. After
	precomposing with the quasi-isomorphism $\tau^{\leq0}Z\to Z$, we may
	assume that $Z\in\Obj(\cate{C}^{\leq0}(\mathcal{A}))$. Next, the
	adjunction in Proposition~\ref{prop:truncation-adjunction} factors $s$
	and $a$ through the quasi-isomorphism
	$Z\twoheadrightarrow\tau^{\geq-n}Z$. Hence we may assume that
	$Z\in\Obj(\cate{C}^{[-n,0]}(\mathcal{A}))$.

	Since $Z$ is exact outside degree $0$, imitate the definition of the
	pushout of an extension above and define a complex $Z'$ to be the
	second row of the following diagram:
% segment-id: o014.li2.chapter4.morphisms-extensions.g012
	\[\begin{tikzcd}[every cell/.append style={font=\small}]
		\cdots 0 \arrow[r] & Z^{-n} \arrow[r] \arrow[d, "{a^{-n}}"'] \arrow[phantom, rd, "\boxplus" description] & Z^{-n+1} \arrow[r] \arrow[d] & Z^{-n+2} \arrow[r] \arrow[equal, d] & \cdots \\
		\cdots 0 \arrow[r] & Y \arrow[r] & Z^{-n+1} \dsqcup{Z^{-n}} Y \arrow[r] & Z^{-n+2} \arrow[r] & \cdots
	\end{tikzcd}\]
	The vertical morphisms give a quasi-isomorphism $Z\to Z'$, and $a$
	factors as $Z\to Z'\xrightarrow{a'}Y[n]$.

	Moreover, $Z\xrightarrow{s}X$ also factors canonically as
	$Z\to Z'\xrightarrow{s'}X$. This is immediate when $n>1$; when $n=1$,
	the morphism $(s')^0:Z^0\dsqcup{Z^{-1}}Y\to X$ is determined by
	$s^0:Z^0\to X$ and $Y\xrightarrow{0}X$.

	In short, we have reduced to the case in which
	$X\xleftarrow{s}Z\xrightarrow{a}Y[n]$ has the form
% segment-id: o014.li2.chapter4.morphisms-extensions.g013
	\[\begin{tikzcd}
		\cdots 0 \arrow[r] & 0 \arrow[r] & 0 \arrow[r] & \cdots \arrow[r] & X \arrow[r] & 0 \cdots \\
		\cdots 0 \arrow[r] & Z^{-n} \arrow[r] \arrow[equal, d] \arrow[u] & Z^{-n+1} \arrow[r] \arrow[u] \arrow[d] & \cdots \arrow[r] & Z^0 \arrow[u] \arrow[d] \arrow[r] & 0 \cdots \\
		\cdots 0 \arrow[r] & Y \arrow[r] & 0 \arrow[r] & \cdots \arrow[r] & 0 \arrow[r] & 0 \cdots
	\end{tikzcd}\]
	Flattening this diagram, however, gives the $n$-extension
	$0\to Y\to Z^{-n+1}\to\cdots\to Z^0\to X\to0$. This proves
	surjectivity.

	We now prove injectivity. In view of \eqref{eqn:localization-sim}, the
	starting point is a commutative diagram in $\cate{K}(\mathcal{A})$
% segment-id: o014.li2.chapter4.morphisms-extensions.g014
	{\scriptsize
	\[\begin{tikzcd}
		& Z_2 \arrow[ld, "s_2"'] \arrow[rd, "a_2"] & \\
		X & W \arrow[l, "t"'] \arrow[r, "b"] \arrow[u] \arrow[d] & Y[n] \\
		& Z_1 \arrow[lu, "s_1"] \arrow[ru, "a_1"'] &
	\end{tikzcd} \quad \begin{array}{l}
		W \in \Obj(\cate{C}(\mathcal{A})), \quad s_i,t:\;\text{quasi-isomorphisms}, \\
		(Z_i;s_i,a_i)\;\text{is a reduced \(n\)-extension},\\
		(i=1,2).
	\end{array}\]
	}
	We wish to show that this diagram gives an equivalence of
	$n$-extensions. As before, first precompose with
	$\tau^{\leq0}W\to W$ to reduce to
	$W\in\Obj(\cate{C}^{\leq0}(\mathcal{A}))$. Next, use the truncation
	adjunction to factor all the morphisms in $\cate{C}(\mathcal{A})$
	through $W\twoheadrightarrow\tau^{\geq-n}W$. Let $K$ be the kernel of
	$W\twoheadrightarrow\tau^{\geq-n}W$. Since
	$\Hom^{-1}(K,V)=0$ for every
	$V\in\Obj(\cate{C}^{[-n,0]}(\mathcal{A}))$, the homotopies between
	these morphisms factor in the same way. This operation therefore does
	not change the commutativity of the diagram in
	$\cate{K}(\mathcal{A})$. We have reduced to the case
	$W\in\Obj(\cate{C}^{[-n,0]}(\mathcal{A}))$.

	Similarly, $\Hom^{-1}(W,X)=0$, so the left half of the diagram already
	commutes in $\cate{C}(\mathcal{A})$. For the right half, since
	$a_i^{-n}=\identity_Y$, we may modify the morphisms $W\to Z_i$ by
	suitable homotopies so that the right half also commutes in
	$\cate{C}(\mathcal{A})$. Now repeat the pushout operation from the
	surjectivity proof to reduce to $W^{-n}=Y$ and
	$b^{-n}=\identity_Y$, without changing the commutativity of the left
	half, as one verifies directly. Thus $(W;t,b)$ also corresponds to an
	$n$-extension, and the commutative diagram witnesses the equivalence
	relation $\sim$ in Definition~\ref{def:Yoneda-Ext}. This proves
	injectivity.

	Compatibility of composition, direct sum, pullback, and pushout with
	this bijection is a routine verification. For the Baer sum $\dot{+}$,
	Proposition~\ref{prop:additive-cat-addition} realizes the sum $f+g$ in
	$\Ext^n(X,Y):=\Hom_{\cate{D}(\mathcal{A})}(X,Y[n])$ as the composite
	$X\to X\oplus X\xrightarrow{f\oplus g}(Y\oplus Y)[n]\to Y[n]$,
	which agrees exactly with the definition of $\dot{+}$.
\end{proof}

% segment-id: o014.li2.chapter4.morphisms-extensions.p018
Dually, an element of $\Hom_{\cate{D}(\mathcal{A})}(X,Y[n])$ can also be
represented by a diagram of the form
$X\xrightarrow{b}Z\xleftarrow{t}Y[n]$ in $\cate{C}(\mathcal{A})$, where
$t$ is a quasi-isomorphism. The exercises for this chapter compare these
two constructions.

% segment-id: o014.li2.chapter4.morphisms-extensions.pr005
\begin{proposition}\label{prop:Ext1-via-ses}
	Consider an extension
	$0\to Y\xrightarrow{f}E\xrightarrow{g}X\to0$. In the exact sequence
	from Proposition~\ref{prop:Ext-generalities}~(iii),
% segment-id: o014.li2.chapter4.morphisms-extensions.q011
	\begin{gather*}
		\Hom_{\mathcal{A}}(X,E) \xrightarrow{g_*} \Hom_{\mathcal{A}}(X,X) \to \Ext^1(X,Y),
	\end{gather*}
	the element of $\Ext^1(X,Y)$ determined by this extension is the image
	of $\identity_X$.
\end{proposition}

% segment-id: o014.li2.chapter4.morphisms-extensions.pf009
\begin{proof}
	The connecting homomorphism
	$\Hom_{\mathcal{A}}(X,X)\to\Ext^1(X,Y)$ is precisely $e_*$, where
	$e\in\Hom_{\cate{D}(\mathcal{A})}(X,Y[1])$ is determined by the diagram
	in $\cate{C}(\mathcal{A})$
% segment-id: o014.li2.chapter4.morphisms-extensions.g015
	\[\begin{tikzcd}
		X & \Cone(f) \arrow[r, "{\beta(f)}"] \arrow[l, "\Phi"', "\text{quasi-isomorphism}" inner sep=0.6em] & {Y[1]}
	\end{tikzcd}\]
	as in Proposition~\ref{prop:ses-vs-triangle}. But $\Cone(f)$ is the
	complex $Y\xrightarrow{f}E$ concentrated in degrees $-1$ and $0$.
	Comparison with Theorem~\ref{prop:Yoneda-Ext} immediately shows that
	$e$ comes from the extension
	$0\to Y\xrightarrow{f}E\xrightarrow{g}X\to0$.
\end{proof}

% segment-id: o014.li2.chapter4.morphisms-extensions.pr006
\begin{proposition}\label{prop:Yoneda-zero}
	The zero element of the additive group
	$\left(\Ext^{n,\text{Yoneda}}(X,Y),\dot{+}\right)$ is represented by
	the following extensions:
% segment-id: o014.li2.chapter4.morphisms-extensions.tb001
	\begin{center}\begin{tabular}{|c|c|} \hline
		$n = 1$ & The split extension $0 \to Y \to X \oplus Y \to X \to 0$ \\ \hline
		$n = 2$ & $0 \to Y \xrightarrow{\identity_Y} Y \xrightarrow{0} X \xrightarrow{\identity_X} X \to 0$ \\ \hline
		$n \geq 3$ & $0 \to Y \xrightarrow{\identity_Y} Y \to 0 \to \cdots \to 0 \to X \xrightarrow{\identity_X} X \to 0$ . \\ \hline
	\end{tabular}\end{center}
\end{proposition}

% segment-id: o014.li2.chapter4.morphisms-extensions.pf010
\begin{proof}
	For $n=1$, suppose that the short exact sequence
	$0\to Y\xrightarrow{f}E\xrightarrow{g}X\to0$ corresponds to the zero
	element. The exact sequence in Proposition~\ref{prop:Ext1-via-ses}
	shows that there is an $s\in\Hom_{\mathcal{A}}(X,E)$ such that
	$gs=\identity_X$; hence the short exact sequence splits.

	Next consider $n=2$. Both $\Ext^1(0,Y)$ and $\Ext^1(X,0)$ are plainly
	zero groups; corresponding extensions may be taken to be
	$0\to Y\xrightarrow{\identity}Y\to0\to0$ and
	$0\to0\to X\xrightarrow{\identity}X\to0$. Since
	$\Ext^1\times\Ext^1\to\Ext^2$ is bilinear, splicing these two
	extensions gives the zero element of $\Ext^2(X,Y)$. For $n\geq3$,
	consider the $(n-1)$-extension
	$0\to\cdots\to0\to X\xrightarrow{\identity}X\to0$; the same argument
	gives the zero element of $\Ext^n(X,Y)$.
\end{proof}
